\uuid{FgtZ}
\exo7id{3761}
\titre{exo7 3761}
\auteur{quercia}
\organisation{exo7}
\datecreate{2010-03-11}
\isIndication{false}
\isCorrection{true}
\chapitre{Espace euclidien, espace normé}
\sousChapitre{Projection, symétrie}
\cpge{mpsi-mp2i.applications-lineaires.b, pcsi.applications-lineaires.b, ptsi.applications-lineaires.b, mp-mpi.endomorphismes-euclidien.b, pc.endomorphismes-euclidiens.b, psi.endomorphismes-euclidien.b, pt.isometries-euclidien.b, mp-mpi.endomorphismes-euclidien.c, pc.endomorphismes-euclidiens.a, psi.endomorphismes-euclidien.a, pt.isometries-euclidien.a, mp-mpi.endomorphismes-euclidien.e, psi.endomorphismes-euclidien.e, pt.isometries-euclidien.d}
\module{Algèbre}
\niveau{L2}
\difficulte{}

\contenu{
\question{
On considère l'endomorphisme de~$\R^3$ de matrice dans la base canonique~:
$A = \begin{pmatrix}a^2 &ab-c &ac+b\cr ab+c &b^2 &bc-a\cr ac-b &bc+a &c^2\cr\end{pmatrix}$
avec~$a,b,c\in\R$.
Déterminer $a,b,c$ de sorte que~$f$ soit une isométrie, et la préciser.
}
\reponse{
$^t\!AA = I \Leftrightarrow a^2+b^2+c^2=1$. Dans ce cas $\det(A) = \mathrm{tr}(A) = 1$ donc $f$ est un quart de tour.
 L'axe du quart de tour est engendré par $\begin{pmatrix}a\cr b\cr c\cr\end{pmatrix}$.
}
}
